% !TEX root = Tesi_Corsi_Leonardo_656276.tex
\chapter{Progetto logico}
\label{ch:progetto}

Il capitolo descrive le scelte di progetto prima di qualunque riga di codice, e
procede dal generale al particolare. 

Le prime due sezioni definiscono il modello dell'array e la cella che lo compone. 

La terza tratta i dettagli di 
progettazione del canale, che è l'unico meccanismo di interazione dell'intero 
sistema e ne concentra tutte le scelte non ovvie. 

La quarta mostra come il programma host interagisca con il perimetro senza
che questo sia un caso speciale per le celle. 

La quinta discute i due piani di parallelismo del
progetto, quello dell'array simulato e quello del simulatore, che il resto
della tesi tiene distinti.

Da qui in poi, il termine \emph{simulatore} indica l'intero programma, ovvero
la parte che fa avanzare l'array ciclo di clock dopo ciclo di clock e il \emph{programma host},
che immette e consuma i valori sul perimetro e ne raccoglie i risultati.
L'\emph{array} e le sue celle esistono soltanto dentro il simulatore.
La distinzione vale per tutte le grandezze che seguono: cicli, istruzioni, canali
e protocollo si riferiscono all'array simulato, thread e speedup al simulatore.

\section{Il modello dell'array}
\label{sec:architettura}

NESO simula una griglia rettangolare di $R \times C$ celle. Ogni cella è un
processore RISC-V completo e indipendente: ha i suoi registri, il suo
\emph{program counter} e la sua memoria privata. Non esiste memoria condivisa,
non esiste un bus, non esiste un nodo di coordinamento. L'unico modo che una cella 
ha per influenzare il resto del sistema è mandare una parola a uno dei quattro vicini 
ortogonali, e l'unico modo di esserne influenzata è riceverne una.

Le regole del modello sono quattro:

\begin{enumerate}
    \item \emph{connessioni solo locali.} La cella in posizione $(r,c)$ parla
          con $(r-1,c)$, $(r,c+1)$, $(r+1,c)$ e $(r,c-1)$, e con nessun altro.
          Le direzioni sono nominate nord, est, sud e ovest e numerate in un
          solo modo in tutto il progetto (Nord = 0, Est = 1, Sud = 2, Ovest = 3), 
          così che la direzione opposta si ottenga invertendo con uno \emph{xor}
          il più significativo dei due bit che la rappresentano:
          $\mathrm{opp}(d) = d \oplus 2$.
    \item \emph{tempo globale e discreto.} Esiste un clock unico. A ogni ciclo
          ogni cella ancora attiva esegue esattamente un'istruzione.
    \item \emph{canali monodirezionali di una posizione (grado di asincronia
          uno).} Fra due celle
          adiacenti esistono due canali, di direzione opposta, e ciascuno può contenere
          al più una parola non ancora letta.
        \item \emph{un solo programma.} Tutte le celle eseguono lo stesso programma,
          di cui ciascuna ha una copia nella propria memoria; il ruolo che
          ciascuna assume lo decide leggendo la propria posizione.

\end{enumerate}

L'ultima regola è una scelta del modello di programmazione, non un vincolo
dell'architettura. La realizzabilità in hardware si regge sull'uniformità
delle celle, non sull'identità dei programmi: le celle hanno tutte la stessa
ISA e la stessa interfaccia verso i vicini, quindi se ne costruisce una e la
si replica $R \times C$ volte. Nulla impedirebbe di caricare in celle diverse
programmi diversi, tutti RV32I: cambierebbe soltanto il \emph{boot} logico
dell'array, cioè la fase in cui il simulatore deposita le immagini nelle
memorie private. Tenere un solo programma semplifica quel caricamento e fa
del ruolo di ciascuna cella una conseguenza della sua posizione.

È il modello di esecuzione SPMD (\emph{Single Program, Multiple
Data})~\cite{Darema}: un solo codice binario gira su tutte le celle, e a differenziarne
il comportamento sono eventuali salti condizionati a seconda della propria posizione
nell'array, che è un'informazione già nota ai nodi per costruzione. 
Nel programma produttore-consumatore, per esempio, la cella di colonna 0
produce e le altre consumano. Il meccanismo con cui una cella conosce la
propria posizione è nella sezione~\ref{sec:registri}. La
figura~\ref{fig:array} mostra la topologia e il cablaggio appena descritti.



\section{La cella}
\label{sec:cella}

\begin{figure}[htbp]
    \centering
    \begin{tikzpicture}[
        scale=0.95, transform shape,
        cella/.style={draw, thick, minimum size=1.5cm, font=\small},
        canale/.style={-{Latex[length=1.8mm]}, shorten >=1pt, shorten <=1pt},
        bordo/.style={canale, gray, dashed}
    ]
    % le nove celle
    \foreach \r in {0,1,2} {
        \foreach \c in {0,1,2} {
            \node[cella] (n\r\c) at (3*\c, -3*\r) {$(\r,\c)$};
        }
    }
    % canali est-ovest: corsia alta verso est, corsia bassa verso ovest
    \foreach \r in {0,1,2} {
        \foreach \c/\cc in {0/1, 1/2} {
            \draw[canale] ([yshift=3mm]n\r\c.east)  -- ([yshift=3mm]n\r\cc.west);
            \draw[canale] ([yshift=-3mm]n\r\cc.west) -- ([yshift=-3mm]n\r\c.east);
        }
    }
    % canali nord-sud: corsia sinistra verso sud, corsia destra verso nord
    \foreach \c in {0,1,2} {
        \foreach \r/\rr in {0/1, 1/2} {
            \draw[canale] ([xshift=-3mm]n\r\c.south)  -- ([xshift=-3mm]n\rr\c.north);
            \draw[canale] ([xshift=3mm]n\rr\c.north) -- ([xshift=3mm]n\r\c.south);
        }
    }
    % il bordo: i canali di perimetro
    \foreach \c in {0,1,2} {
        \draw[bordo] ([xshift=-3mm]$(n0\c.north)+(0,0.8)$) -- ([xshift=-3mm]n0\c.north);
        \draw[bordo] ([xshift=3mm]n0\c.north) -- ([xshift=3mm]$(n0\c.north)+(0,0.8)$);
        \draw[bordo] ([xshift=-3mm]n2\c.south) -- ([xshift=-3mm]$(n2\c.south)-(0,0.8)$);
        \draw[bordo] ([xshift=3mm]$(n2\c.south)-(0,0.8)$) -- ([xshift=3mm]n2\c.south);
    }
    \foreach \r in {0,1,2} {
        \draw[bordo] ([yshift=3mm]$(n\r0.west)-(0.8,0)$) -- ([yshift=3mm]n\r0.west);
        \draw[bordo] ([yshift=-3mm]n\r0.west) -- ([yshift=-3mm]$(n\r0.west)-(0.8,0)$);
        \draw[bordo] ([yshift=3mm]n\r2.east) -- ([yshift=3mm]$(n\r2.east)+(0.8,0)$);
        \draw[bordo] ([yshift=-3mm]$(n\r2.east)+(0.8,0)$) -- ([yshift=-3mm]n\r2.east);
    }
    % tutto cio' che sta fuori dall'array e' l'host
    \draw[gray, dashed, rounded corners]
        ($(n00.north west)+(-1.3,1.3)$) rectangle ($(n22.south east)+(1.3,-1.3)$);
    \node[gray, font=\small\itshape, fill=white, inner sep=2pt]
        at ($(n02.north east)+(0.35,1.3)$) {host};
    \end{tikzpicture}
    \caption[Topologia dell'array e canali verso i quattro vicini]{Topologia
    dell'array. Fra due celle adiacenti corrono due canali monodirezionali, uno
    per verso. Sul perimetro le direzioni che punterebbero fuori dalla griglia
    sono collegate a canali di bordo (tratteggiati), che l'host
    alimenta e consuma con le stesse operazioni di una cella vicina. La coppia
    $(r,c)$ che etichetta ogni cella è la sua posizione nell'array, riga e
    colonna, sulla quale il programma si dirama per differenziare il
    comportamento dei singoli nodi (sezione~\ref{sec:registri}).}
    \label{fig:array}
\end{figure}

Una cella contiene lo stato architetturale di un processore RV32I, esteso ora
con i canali con cui comunica con i vicini, e la memoria privata su cui lavora.

Lo stato di una cella comprende quindi:

\begin{itemize}
    \item i trentadue registri interi dell'architettura, con
          \texttt{x0} cablato a zero;
    \item il \emph{program counter};
    \item quattro canali di uscita, uno per direzione, di cui la cella è
          l'unica proprietaria;
    \item quattro riferimenti a canali di ingresso, che sono i canali di
          uscita dei vicini;
    \item due contatori di misura, tenuti dal simulatore a supporto delle misure
          descritte nella sezione~\ref{sec:metodologia};
    \item una memoria privata di 16 KB, indirizzata a byte, che contiene sia il 
          codice sia i dati.
\end{itemize}
La memoria privata è il punto in cui il simulatore risulta più vincolante del necessario:
16\,KB per cella significa che una griglia $12 \times 12$ occupa
già più di due megabyte di sola memoria simulata. La capacità della memoria privata è però una
costante che può essere modificata a tempo di compilazione (\texttt{MEM\_SIZE} in 
\texttt{risc.h}), non un vincolo intrinseco del modello: per griglie più grandi si può 
ridurre tale costante e ricompilare senza necessità di modificare l'architettura. La
figura~\ref{fig:cella} riassume gli elementi appena elencati.

\begin{figure}[htbp]
    \centering
    \begin{tikzpicture}[
        scale=0.8, transform shape,
        blocco/.style={draw, thick, minimum height=0.9cm, font=\small, align=center},
        ch/.style={draw, fill=black!8, minimum width=1.5cm, minimum height=0.55cm,
                   font=\scriptsize\ttfamily},
        proprio/.style={-{Latex[length=1.8mm]}},
        altrui/.style={-{Latex[length=1.8mm]}, gray, dashed},
        eti/.style={font=\scriptsize, gray}
    ]
    % il corpo della cella
    \node[draw, thick, rounded corners, minimum width=8.4cm, minimum height=3.6cm]
        (cella) {};
    \node[font=\small\itshape, anchor=north west]
        at ([shift={(0.15,-0.1)}]cella.north west) {cella $(r,c)$};

    \node[blocco, minimum width=3.4cm] (regs)
        at ([shift={(-1.9,0.35)}]cella.center) {registri \texttt{x0}\dots\texttt{x31}};
    \node[blocco, minimum width=2.6cm] (pc)
        at ([shift={(2.1,0.35)}]cella.center) {\texttt{pc}};
    \node[blocco, minimum width=7.4cm] (mem)
        at ([shift={(0,-0.9)}]cella.center) {memoria privata: 16\,KB};

    % i quattro canali di uscita: la cella ne e' l'unica proprietaria
    \node[ch] (oN) at ([shift={(-1.6,1.5)}]cella.north) {out\_ch[N]};
    \node[ch] (oS) at ([shift={(-1.6,-1.5)}]cella.south) {out\_ch[S]};
    \node[ch] (oO) at ([shift={(-1.9,0.7)}]cella.west)  {out\_ch[O]};
    \node[ch] (oE) at ([shift={(1.9,0.7)}]cella.east)   {out\_ch[E]};

    \draw[proprio] ([xshift=-1.6cm]cella.north) -- (oN);
    \draw[proprio] ([xshift=-1.6cm]cella.south) -- (oS);
    \draw[proprio] ([yshift=0.7cm]cella.west)   -- (oO);
    \draw[proprio] ([yshift=0.7cm]cella.east)   -- (oE);

    % i quattro ingressi: sono i canali di uscita dei vicini, non sono suoi
    \draw[altrui] ([shift={(1.6,1.5)}]cella.north) -- ([xshift=1.6cm]cella.north);
    \draw[altrui] ([shift={(1.6,-1.5)}]cella.south) -- ([xshift=1.6cm]cella.south);
    \draw[altrui] ([shift={(-1.9,-0.7)}]cella.west) -- ([yshift=-0.7cm]cella.west);
    \draw[altrui] ([shift={(1.9,-0.7)}]cella.east)  -- ([yshift=-0.7cm]cella.east);

    \node[eti, anchor=south] at ([shift={(1.6,1.55)}]cella.north) {\texttt{in\_ch[N]}};
    \node[eti, anchor=north] at ([shift={(1.6,-1.55)}]cella.south) {\texttt{in\_ch[S]}};
    \node[eti, anchor=east]  at ([shift={(-1.95,-0.7)}]cella.west) {\texttt{in\_ch[O]}};
    \node[eti, anchor=west]  at ([shift={(1.95,-0.7)}]cella.east)  {\texttt{in\_ch[E]}};
    \end{tikzpicture}
    \caption[Struttura di una cella]{Struttura di una cella. I quattro
    canali di uscita appartengono alla cella, che ne è l'unica proprietaria; i
    quattro ingressi (tratteggiati) non sono suoi, sono i canali di uscita dei
    quattro vicini, raggiunti per puntatore. La memoria è privata della cella, ovvero 
    non è visibile a nessun'altra.}
    \label{fig:cella}
\end{figure}


\section{Il canale}
\label{sec:canale}

Il canale è il componente su cui si basa il protocollo di comunicazione, ed è
deliberatamente molto compatto.

\begin{lstlisting}[caption={La struttura del canale e i suoi due test},label={lst:channel},language=C,float=htbp]
typedef struct Channel {
    uint32_t data;       /* visibile al consumatore */
    uint8_t  wp, rp;     /* i due ready bit, mod 2 */

    uint32_t data_next;  /* registro di uscita: lo carica OUT */
    uint8_t  wp_next, rp_next;
} Channel;

/* i test guardano lo stato ATTUALE, mai i gemelli _next */
static inline int ch_isrdy(const Channel *c) {   /* c'è una parola da leggere? */
    return c -> wp != c -> rp;
}
static inline int ch_iswrt(const Channel *c) {   /* si può pubblicare? */
    return c -> wp == c -> rp;
}

/* IN: legge e, se c'era un dato, libera la posizione */
static inline uint32_t ch_read_c(Channel *c) {
    uint32_t v = c -> data;
    if (ch_isrdy(c)) {
        c -> rp_next = c -> rp ^ 1;
    }
    return v;
}
\end{lstlisting}
\clearpage
Tutto lo stato necessario a implementare il protocollo è costituito da due 
contatori modulo 2 (a un bit) e dal comparatore che li
confronta. Nel codice si chiamano \texttt{wp} e \texttt{rp}, dal
\emph{write pointer} e \emph{read pointer} da cui derivano, ma nel seguito
saranno i due \emph{ready bit} del canale, uno per lato: il produttore commuta
il proprio per dire che c'è un valore nuovo, il consumatore commuta il proprio
per dire che l'ha preso. Nessun campo registra l'occupazione del canale, e lo 
stato si legge confrontando i due contatori invece che il valore dell'uno o dell'altro.
Ogni campo ha poi un gemello con il suffisso \texttt{\_next}, che l'aggiornamento di fine ciclo usa
per tenere separato ciò che un ciclo legge da ciò che vi scrive (sezione~\ref{sec:due-fasi}).

I due test del listato registrano proprio quel confronto, e va notato quali
campi guardano: \texttt{wp} e \texttt{rp}, cioè lo stato \emph{attuale}, mai i
gemelli \texttt{\_next}. È un dettaglio da cui dipende l'intero protocollo, e
la sezione~\ref{sec:finestra} ne mostra la ragione. Il confronto fra i due
campi si riassume così:

\begin{center}
\begin{tabular}{ll}
    \toprule
    il consumatore ha una parola da leggere & $\mathtt{wp} \neq \mathtt{rp}$ \\
    il produttore può pubblicare & $\mathtt{wp} = \mathtt{rp}$ \\
    \bottomrule
\end{tabular}
\end{center}

La scelta è fra due famiglie di protocolli di sincronizzazione: uno
\emph{a livelli}, in cui lo stato del canale è il valore assoluto di un
segnale condiviso, e uno \emph{a transizione di livello}, in cui è la commutazione
di un segnale, non il suo valore, a portare l'informazione~\cite{Sutherland1989}.
Il canale di NESO adotta un protocollo della seconda famiglia.

Un unico bit di occupazione non basterebbe, perché andrebbe scritto da due lati
diversi: il produttore lo setterebbe pubblicando, il consumatore lo resetterebbe
leggendo. Le due cose accadono nello stesso ciclo ogni volta che il vicino
consuma mentre il mittente pubblica, e le due scritture finirebbero sullo stesso
campo. Nel simulatore l'esito dipenderebbe dall'ordine con cui le celle vengono
visitate; in hardware non si può realizzare un circuito in cui due sorgenti pilotano lo stesso flip-flop nel
medesimo istante.

Con due contatori, ogni lato scrive soltanto il proprio campo: il produttore
tocca solo \texttt{wp}, il consumatore solo \texttt{rp}. Non c'è contesa e non
c'è dipendenza dall'ordine di scansione, perché nessuno dei due scrive una
variabile comune. I contatori sono \emph{modulo due} perché il canale ha una
sola posizione: bastano due stati, quindi un solo bit, che per costruzione non
può andare in \emph{overflow}.

Ogni canale è condiviso dalle due celle che collega: il canale di uscita verso
est del mittente \emph{è} il canale di ingresso da ovest del ricevente. Non
esiste un livello di rete e non esiste una tabella di instradamento. Nel
simulatore il cablaggio della griglia realizza questa condivisione assegnando
puntatori, senza copiare nulla.


\section{Il bordo}
\label{sec:bordo}

Le celle del perimetro hanno direzioni che puntano fuori dalla griglia. Il
simulatore alloca allora $2(R+C)$ canali di bordo e li collega a quelle
direzioni, così che ogni cella abbia sempre quattro ingressi e quattro uscite
validi. Da quel momento l'host può fare ciò che farebbe un vicino, con lo stesso
protocollo della comunicazione interna: caricare un valore e pubblicarlo, oppure
verificare che ci sia un dato e consumarlo.

La conseguenza è che nessuna cella ha la percezione di stare sul bordo. Il programma
che gira nell'angolo è lo stesso che gira al centro, legge i quattro ingressi
allo stesso modo e pubblica sulle quattro uscite allo stesso modo. Questa
astrazione rende possibile il programma completamente uniforme discusso
nella sezione~\ref{sec:kernel-jacobi}.


\section{Parallelizzazione}
\label{sec:parallelizzazione}

Il progetto è parallelo su due piani distinti, che rispondono a domande diverse.

\subsection{Il calcolo distribuito sulle celle}

Il primo piano è quello dell'array: a ogni ciclo ogni cella attiva calcola il
proprio nuovo stato a partire dallo stato in cui l'array si trova all'inizio del
ciclo, che è lo stesso per tutte.
La forma che il parallelismo assume qui è quella sistolica: i dati si
muovono e il calcolo sta fermo. Ogni cella può accumulare un risultato che non
la lascia mai, e a scorrere sono gli operandi che la attraversano. Il programma
di moltiplicazione di matrici, descritto nella sezione~\ref{sec:kernel}, ne è
il caso canonico.

Poiché tutte le celle eseguono lo stesso programma, la suddivisione
del lavoro non è espressa da un programma di coordinamento ma dalla posizione:
ogni cella nasce sapendo dove si trova e sceglie il proprio ruolo da lì
(sezione~\ref{sec:registri}). Il perimetro ha $O(R+C)$ canali contro gli
$O(R \cdot C)$ della griglia, e i dati dall'esterno entrano ed escono solo di
lì: al crescere della griglia la quota di comunicazione a carico dell'host
diminuisce, e la sezione~\ref{sec:scalabilita} la misura.

Nell'implementazione attuale quell'ingresso e quell'uscita costano $O(R+C)$
anche in tempo, perché \texttt{grid\_border\_fill} e \texttt{grid\_border\_drain}
(sezione~\ref{sec:io-bordo}) sono cicli seriali sull'host: non è un limite del
modello, dato che ogni canale di bordo è indipendente dagli altri e vale lo
stesso argomento dell'unico scrittore per campo; è un costo del simulatore che
una versione parallela di quei due cicli eliminerebbe.

\subsection{La simulazione di un ciclo di clock}

Il secondo piano è quello del simulatore, che per far avanzare l'array di un
ciclo di clock può usare più processori della macchina su cui gira. Il codice
non ha avuto bisogno di alcun \emph{lock}, per una ragione strutturale: il ciclo
di clock separa una fase di calcolo da una fase di aggiornamento
(sezione~\ref{sec:due-fasi}), e in fase di calcolo ogni campo ha un solo
scrittore:

\begin{center}
\begin{tabular}{ll}
    \toprule
    campo & unico scrittore \\
    \midrule
    \texttt{data\_next}, \texttt{wp\_next} & il produttore del canale \\
    \texttt{rp\_next} & il consumatore, unico per costruzione \\
    registri, PC, memoria, contatori & la cella proprietaria \\
    \bottomrule
\end{tabular}
\end{center}

Ciascuna delle due fasi è quindi un ciclo banalmente parallelo (per usare la
traduzione consueta di \emph{embarrassingly parallel}): le iterazioni sono
indipendenti e non richiedono sincronizzazione al proprio interno. 
Le uniche sincronizzazioni necessarie sono le due barriere di ogni ciclo, fra
calcolo e aggiornamento e fra aggiornamento e ciclo successivo, discusse nella
sezione~\ref{sec:omp-impl}.

Un caso che potrebbe sembrare una \emph{race condition} in realtà non lo è: una cella
che legge da un vicino scrive \texttt{rp\_next} \emph{dentro la struttura di
quel vicino}. Due thread lavorano quindi sulla stessa struttura, ma su campi
diversi, e in C11 membri distinti di una struttura sono posizioni di memoria
distinte~\cite{C11}: non c'è \emph{race condition} nemmeno fra i due byte adiacenti. È una
questione diversa dal \emph{false sharing}\footnote{Due thread che scrivono
campi distinti ma vicini si contendono ugualmente la linea di cache che li
contiene, perché la coerenza fra le cache si mantiene a granularità di linea e
non di variabile: ogni scrittura invalida la copia dell'altro, che deve
rileggerla. È condivisione solo apparente, da cui il nome, e a differenza di una
\emph{race condition} non altera il risultato, lo rallenta soltanto~\cite{Torrellas}.},
che riguarda le prestazioni e non la correttezza, e di cui si parla nella
sezione~\ref{sec:prestazioni}.

La separazione in due fasi serve al determinismo del risultato. Che renda
superfluo ogni \emph{lock} nella simulazione di un singolo ciclo di clock è una
conseguenza dell'unico scrittore per campo.
